\documentclass[10pt]{article}
\usepackage[margin=0.72in]{geometry}
\usepackage[T1]{fontenc}
\usepackage[utf8]{inputenc}
\usepackage[table]{xcolor}
\usepackage{booktabs}
\usepackage{array}
\usepackage{tabularx}
\usepackage[hidelinks]{hyperref}
\setlength{\parindent}{0pt}
\setlength{\parskip}{0.40em}
\setlength{\emergencystretch}{2em}

\newcommand{\revcomment}[1]{\textcolor{black}{\textbf{Reviewer comment. #1}}}
\newcommand{\questionsep}{\par\vspace{1.15em}}
\newcommand{\response}{\textbf{Response.}\enspace}
\definecolor{inhouseblue}{RGB}{225,240,255}

\begin{document}

\begin{center}
{\Large Response to Reviewer aSvP}\\[2pt]
Submission 13175
\end{center}

\textbf{Opening response.}
We sincerely thank the reviewer for recognizing the potential value of sMMC and
for identifying the central question we had not answered clearly enough: what
research becomes possible because the resource is larger, cell aligned, and
context rich? This resource has accompanied me since my first PhD year, when I
was new to spatial transcriptomics. As I approach graduation, I hope its unified
release helps new students enter the field with fewer practical barriers.

We also clarify the new primary component. The current working manifest contains
99 unique samples (100 processing records) and 28,315,247 aligned targets. Our
in-house DBiC-seq collection contributes 21 samples and 54,304 measured cells;
paired-record QC removes 315, leaving 53,989 aligned cells. These records pair
cell morphology, RNA, and cellular context. A broader continuing DBiC-seq pool
now contains approximately 200,000 paired cells but is not counted in the table
below.

\begin{center}
\begingroup
\footnotesize
\setlength{\tabcolsep}{3pt}
\begin{tabularx}{\linewidth}{@{}
>{\raggedright\arraybackslash}p{0.27\linewidth}
>{\raggedright\arraybackslash}p{0.18\linewidth}
>{\raggedright\arraybackslash}Xr@{}}
\toprule
Component & Samples & Platforms & Aligned cells after QC \\
\midrule
\rowcolor{inhouseblue}
\textbf{New in-house primary data} & \textbf{21} & \textbf{DBiC-seq} &
\textbf{53,989} \\
\textbf{Total sMMC-28M working manifest} &
\textbf{99 unique (100 records)} & \textbf{Xenium; Visium HD; DBiC-seq} &
\textbf{28,315,247} \\
\bottomrule
\end{tabularx}
\endgroup
\end{center}

\questionsep
\textbf{aSvP-1.}\quad
\revcomment{The extent to which this new data can truly promote relevant
research progress has not been fully explained.}

\response
We agree. The resource is valuable not because one predictor has solved
histology-to-expression inference, but because it makes three previously
difficult evaluations measurable: scaling under a fixed task,
leakage-controlled breadth across organs, and failures caused by aggregation or
context shift.

\textbf{Scale.}
On native-Xenium sample HHDX011, we fixed the spatial test regions,
training-selected top-50 HVGs, and regression protocol while varying training
cells for seven frozen encoders:

\begin{center}
\begingroup
\small
\begin{tabular}{lrrrr}
\toprule
Training fraction & 5\% & 10\% & 25\% & 100\% \\
\midrule
Seven-encoder mean Gene Pearson & 0.177 & 0.205 & 0.241 & 0.278 \\
\bottomrule
\end{tabular}
\endgroup
\end{center}

Every encoder improved monotonically. This shows that additional cells are
useful within this controlled range; we do not claim a universal scaling law.

\textbf{Breadth and biological calibration.}
We completed a leakage-controlled benchmark across 30 native-Xenium samples, 17
organ labels, two species, and 360,000 spatially stratified cells. Each sample
used four contiguous spatial holdouts and a 5\% train--test buffer; genes were
selected from training data only.

\begin{center}
\begingroup
\scriptsize
\setlength{\tabcolsep}{3pt}
\begin{tabular*}{\linewidth}{@{\extracolsep{\fill}}lrrrrr@{}}
\toprule
Model & Gene P Top-50 & Gene P Top-200 & Gene S Top-50 & Cell P Top-50 &
F1 Top-50 \\
\midrule
\textbf{Image} & \textbf{0.324} & \textbf{0.202} & \textbf{0.291} &
\textbf{0.442} & \textbf{0.413} \\
Coordinate only & 0.046 & 0.031 & 0.041 & 0.288 & 0.283 \\
Spatial KNN & 0.053 & 0.029 & 0.051 & 0.268 & 0.320 \\
Training mean & N/A & N/A & N/A & 0.305 & 0.253 \\
\bottomrule
\end{tabular*}
\endgroup
\end{center}

The top-50 image Gene Pearson has a biological-sample bootstrap 95\% CI of
0.277--0.368 and exceeds coordinate and KNN controls in 28/30 samples. On
marker/HVG overlap, image Gene Pearson is 0.201 versus 0.031 for coordinates and
0.028 for KNN. Cell-type-stratified pseudobulk RMSE is 0.120 for image prediction
versus 0.224/0.187/0.188 for coordinates/KNN/training mean. Because cell-type
labels are expression derived, this is aggregate biological calibration, not
independent ground truth.

\textbf{Resolution.}
On six native-Xenium samples, matched cell, 8-, 16-, and 55-\(\mu\)m supervision
gives Gene Pearson 0.365, 0.365, 0.363, and 0.330. In dense lung, 55.5\%--66.4\% of
55-\(\mu\)m pseudo-spots mix predicted cell types and contain 73.8\%--81.0\% of
evaluated cells, whereas sparse heart shows only 3.5\%--4.1\% mixed spots and
7.0\%--8.3\% affected cells. Thus, cell-level evaluation exposes
density-dependent mixing hidden by coarse averaging; it is not claimed to be
universally easier.

Together, these experiments turn scale, resolution, and context from descriptive
properties into testable variables: users can study data scaling, organ/platform
transfer, aggregation-induced oversmoothing, spatial leakage, and
context-specific failure.

\questionsep
\textbf{aSvP-2.}\quad
\revcomment{The increase of data scale, resolution, and contexts typically
associates with increased heterogeneity, which may pose critical challenges for
data analysis. Please clarify the heterogeneity and how it may influence
downstream analysis.}

\response
We agree and will make heterogeneity an explicit benchmark variable rather than
a hidden nuisance.

\begin{center}
\begingroup
\scriptsize
\renewcommand{\arraystretch}{1.10}
\begin{tabularx}{\linewidth}{@{}p{0.15\linewidth}p{0.35\linewidth}X@{}}
\toprule
Dimension & Observed heterogeneity & Downstream consequence \\
\midrule
Organ size & 26,366--3,559,793 aligned cells (135-fold) &
Cell-weighted pooling can be dominated by large organs. \\
Gene coverage & 450--72,302 genes &
Full-panel metrics are not directly comparable across panels. \\
Platform & Only 10/25 public tissue strata occur on both platforms &
Platform and biology can be confounded. \\
Spatial density & 55-\(\mu\)m mixing affects 7.0\%--8.3\% of cells in sparse
heart but 73.8\%--81.0\% in dense lung &
Spot averaging hides tissue-dependent heterogeneity. \\
Sample transfer & Across 18 same-organ targets, top-50-HVG Gene Pearson averages
0.170 but ranges 0.016--0.372 &
Mean performance can conceal severe target-sample failures. \\
Context & Ovarian AK (60+) \(\to\) AD (60+) versus AK \(\to\) AL (40-)
changes F1 0.289 \(\to\) 0.072 &
Pooled scores can hide subgroup collapse. \\
\bottomrule
\end{tabularx}
\endgroup
\end{center}

The ovarian comparison is age associated but sample/patient confounded, not
causal. Incomplete donor identity also means the 18-target result is
sample-held-out, not patient-held-out. These limitations illustrate why context
and provenance are necessary.

We will therefore: (i) report organ-macro and platform-separated results rather
than only cell-weighted averages; (ii) distinguish spatial in-domain,
sample-held-out, and verified patient-held-out protocols; (iii) report
Average/Worst/Gap together with sample support; (iv) disclose gene-panel overlap
and metadata missingness; and (v) label every field as directly measured,
processed, generated, or model output. We will not infer missing sensitive
attributes.

These revisions also clarify the scientific value: sMMC does not remove
heterogeneity, but makes its downstream effects visible, measurable, and
reproducible across organs, platforms, spatial resolutions, and biological
contexts.

\end{document}
